\documentclass[12pt,a4paper]{article}

\usepackage[indonesian]{babel}
\usepackage[T1]{fontenc}
\usepackage[utf8]{inputenc}
\usepackage{amsmath,amssymb,amsthm,mathtools}
\usepackage{mathrsfs}
\usepackage{geometry}
\usepackage{setspace}
\usepackage{enumitem}
\usepackage{hyperref}
\usepackage{titlesec}
\usepackage{array}

\geometry{
  left=30mm,
  right=25mm,
  top=25mm,
  bottom=30mm
}

\onehalfspacing
\setlength{\parindent}{1.2cm}
\setlength{\parskip}{0.3em}
\hypersetup{colorlinks=true,linkcolor=blue,urlcolor=blue,citecolor=blue}

\titleformat{\section}{\large\bfseries}{\thesection.}{0.6em}{}
\titleformat{\subsection}{\normalsize\bfseries}{\thesubsection.}{0.6em}{}

\theoremstyle{definition}
\newtheorem{definisi}{Definisi}[section]
\newtheorem{contoh}[definisi]{Contoh}
\newtheorem{remark}[definisi]{Catatan}
\newtheorem{latihan}[definisi]{Latihan}

\theoremstyle{plain}
\newtheorem{teorema}[definisi]{Teorema}
\newtheorem{proposisi}[definisi]{Proposisi}
\newtheorem{lema}[definisi]{Lema}
\newtheorem{akibat}[definisi]{Akibat}

\newcommand{\E}{\mathbb{E}}
\newcommand{\PP}{\mathbb{P}}
\newcommand{\R}{\mathbb{R}}
\newcommand{\N}{\mathbb{N}}
\newcommand{\F}{\mathcal{F}}
\newcommand{\B}{\mathcal{B}}
\newcommand{\ind}{\mathbf{1}}

\begin{document}

\begin{center}
  {\Large \textbf{LAPORAN KALKULUS IT\^O}}\\[0.3cm]
  {\large Mata Kuliah Kalkulus Stokastik}\\[0.5cm]
  Teosofi Hidayah Agung\\
  Departemen Matematika ITS\\
  Juni 2026
\end{center}

\section{Pendahuluan}

\subsection{Latar Belakang}

Kalkulus klasik dibangun di atas fungsi-fungsi yang cukup teratur, misalnya
kontinu, terdiferensialkan, atau setidaknya memiliki variasi terbatas.
Dalam banyak persoalan nyata, terutama pada pemodelan keuangan, fisika, dan
biologi, besaran yang diamati justru berubah secara acak terhadap waktu.
Contoh paling penting adalah proses Wiener $W(t)$ atau gerak Brown, yang
memiliki lintasan kontinu hampir pasti tetapi tidak terdiferensialkan di titik
manapun.

Ketidakteraturan lintasan proses Wiener menyebabkan integral Riemann atau
Riemann--Stieltjes klasik tidak dapat langsung dipakai. Pada kalkulus biasa,
untuk integral
\[
  \int_0^t H(s)\,dA(s),
\]
integrator $A$ umumnya diminta memiliki variasi terbatas. Akan tetapi, dengan
probabilitas satu, lintasan proses Wiener tidak mempunyai variasi terbatas.
Karena itu, diperlukan kerangka baru yang dapat menjelaskan arti integral
terhadap $dW(t)$ dan sekaligus menghasilkan aturan perubahan variabel yang
konsisten. Kerangka inilah yang dikenal sebagai kalkulus It\^o.

Selain penting secara teoretis, kalkulus It\^o juga merupakan dasar bagi
persamaan diferensial stokastik
\[
  dX(t)=\mu(t)\,dt+\sigma(t)\,dW(t),
\]
yang dipakai untuk memodelkan sistem dengan komponen deterministik dan acak
secara simultan.

\subsection{Sejarah Singkat}

Perkembangan kalkulus It\^o tidak lepas dari sejarah gerak Brown. Pada tahun
1900, Louis Bachelier menggunakan gerak acak untuk memodelkan harga aset.
Pada tahun 1905, Albert Einstein menjelaskan gerak partikel kecil dalam fluida
melalui model Brownian motion. Pada tahun 1923, Norbert Wiener memberikan
konstruksi matematis yang ketat bagi Brownian motion, sehingga proses ini juga
disebut proses Wiener. Selanjutnya pada dekade 1940-an, Kiyosi It\^o
mengembangkan integral stokastik dan aturan perubahan variabel yang sekarang
menjadi fondasi kalkulus stokastik modern.

\section{Landasan Teori}

\subsection{Proses Wiener}

Dalam laporan ini, $(\Omega,\F,\PP)$ menyatakan ruang probabilitas lengkap
dengan filtrasi $\{\F_t\}_{t\ge 0}$. Proses stokastik $\{W(t)\}_{t\ge 0}$
disebut \textbf{proses Wiener} apabila memenuhi sifat-sifat berikut.
\begin{enumerate}[label=\alph*.]
  \item $W(0)=0$ hampir pasti.
  \item Untuk setiap $0\le s<t$, inkremen $W(t)-W(s)$ berdistribusi normal
        dengan rataan $0$ dan variansi $t-s$.
  \item Inkremen pada interval-interval saling lepas adalah independen.
  \item Lintasan $t\mapsto W(t,\omega)$ kontinu hampir pasti.
\end{enumerate}

Sifat kedua dan ketiga sangat penting dalam hampir semua pembuktian pada
kalkulus It\^o. Secara khusus, untuk $\Delta_jW:=W(t_{j+1})-W(t_j)$ berlaku
\[
  \E(\Delta_jW)=0,
  \qquad
  \E((\Delta_jW)^2)=t_{j+1}-t_j.
\]

\subsection{Ekspektasi Bersyarat}

Karena pembuktian integral It\^o banyak menggunakan ekspektasi bersyarat,
beberapa sifat yang akan dipakai dirangkum berikut.

\begin{proposisi}
  Misalkan $\mathcal{G}\subseteq\F$ sub-$\sigma$-aljabar dan $\xi,\zeta$ adalah
  peubah acak terintegralkan. Maka berlaku:
  \begin{enumerate}[label=\arabic*.]
    \item \textbf{Linearitas:}
          \[
            \E(\alpha\xi+\beta\zeta\mid\mathcal{G})
            =
            \alpha\E(\xi\mid\mathcal{G})+\beta\E(\zeta\mid\mathcal{G}).
          \]
    \item \textbf{Sifat menara:} jika $\mathcal{H}\subseteq\mathcal{G}$, maka
          \[
            \E(\E(\xi\mid\mathcal{G})\mid\mathcal{H})=\E(\xi\mid\mathcal{H}).
          \]
    \item \textbf{Mengeluarkan faktor terukur:} jika $\xi$ $\mathcal{G}$-terukur
          dan $\xi\zeta$ terintegralkan, maka
          \[
            \E(\xi\zeta\mid\mathcal{G})=\xi\,\E(\zeta\mid\mathcal{G}).
          \]
    \item \textbf{Independensi:} jika $\xi$ independen terhadap $\mathcal{G}$,
          maka
          \[
            \E(\xi\mid\mathcal{G})=\E(\xi).
          \]
    \item \textbf{Mengambil ekspektasi penuh:}
          \[
            \E(\E(\xi\mid\mathcal{G}))=\E(\xi).
          \]
  \end{enumerate}
\end{proposisi}

\subsection{Motivasi Matematis}

Dalam integral Riemann--Stieltjes klasik, konvergensi
\[
  \sum_{j=0}^{n-1} H(s_j)\bigl(A(t_{j+1})-A(t_j)\bigr)
\]
bergantung pada keteraturan fungsi $A$. Untuk proses Wiener, situasinya
berbeda karena lintasannya sangat kasar. Bahkan secara informal,
\[
  \sum_{j=0}^{n-1}|W(t_{j+1})-W(t_j)|
  \longrightarrow \infty
\]
saat ukuran partisi menuju nol. Akibatnya, pendekatan integral klasik tidak
bisa langsung dipakai.

Dalam konteks stokastik, jumlahan
\[
  \sum_{j=0}^{n-1} f(s_j)\bigl(W(t_{j+1})-W(t_j)\bigr)
\]
juga sensitif terhadap pilihan titik $s_j$. Pada integral It\^o, titik yang
dipilih adalah titik kiri $s_j=t_j$. Pilihan ini memastikan integrand hanya
menggunakan informasi yang tersedia sampai waktu $t_j$, sehingga sesuai dengan
gagasan non-antisipatif dalam pemodelan stokastik.

\section{Integral Stokastik It\^o untuk Proses Sederhana}

\subsection{Proses Sederhana Acak}

\begin{definisi}
  Proses stokastik $f(t)$ disebut \textbf{proses sederhana acak} jika terdapat
  partisi berhingga
  \[
    0=t_0<t_1<\cdots<t_n
  \]
  dan peubah acak $\eta_0,\eta_1,\ldots,\eta_{n-1}$ sedemikian sehingga
  \[
    f(t)=\sum_{j=0}^{n-1}\eta_j\ind_{[t_j,t_{j+1})}(t),
  \]
  dengan setiap $\eta_j$ adalah $\F_{t_j}$-terukur dan
  $\E(|\eta_j|^2)<\infty$. Himpunan semua proses sederhana acak dilambangkan
  dengan $M_{\mathrm{step}}^2$.
\end{definisi}

\begin{contoh}
  Pada interval $[0,2]$, definisikan
  \[
    f(t)=2\ind_{[0,1)}(t)+W(1)\ind_{[1,2)}(t).
  \]
  Proses ini merupakan proses sederhana acak karena:
  \begin{enumerate}[label=\alph*.]
    \item koefisien pertama $\eta_0=2$ adalah peubah acak konstan, sehingga
          $\F_0$-terukur;
    \item koefisien kedua $\eta_1=W(1)$ adalah $\F_1$-terukur;
    \item $\E(|\eta_0|^2)=4<\infty$ dan $\E(|\eta_1|^2)=\E(W(1)^2)=1<\infty$.
  \end{enumerate}
  Dengan demikian, semua syarat dalam definisi terpenuhi.
\end{contoh}

\subsection{Definisi Integral It\^o pada Proses Sederhana}

\begin{definisi}
  Jika
  \[
    f(t)=\sum_{j=0}^{n-1}\eta_j\ind_{[t_j,t_{j+1})}(t)
    \in M_{\mathrm{step}}^2,
  \]
  maka \textbf{integral stokastik It\^o} dari $f$ didefinisikan oleh
  \[
    I^\ast(f)
    :=
    \sum_{j=0}^{n-1}\eta_j\bigl(W(t_{j+1})-W(t_j)\bigr).
  \]
\end{definisi}

\begin{contoh}
  Misalkan
  \[
    f(t)=3\ind_{[0,1)}(t)+W(1)\ind_{[1,2)}(t).
  \]
  Berdasarkan definisi, diperoleh
  \[
    I^\ast(f)
    =
    3\bigl(W(1)-W(0)\bigr)+W(1)\bigl(W(2)-W(1)\bigr).
  \]
  Karena $W(0)=0$, maka
  \[
    I^\ast(f)=3W(1)+W(1)\bigl(W(2)-W(1)\bigr).
  \]
  Contoh ini menunjukkan bahwa integral It\^o suatu proses sederhana acak
  menghasilkan peubah acak yang bergantung pada lintasan proses Wiener.
\end{contoh}

\subsection{Sifat Dasar Integral It\^o}

\begin{proposisi}[Isometri It\^o untuk proses sederhana]
  Jika $f\in M_{\mathrm{step}}^2$, maka $I^\ast(f)\in L^2$ dan
  \[
    \E\bigl((I^\ast(f))^2\bigr)
    =
    \E\left(\int_0^\infty f(t)^2\,dt\right).
  \]
\end{proposisi}

\begin{proof}
  Misalkan
  \[
    f(t)=\sum_{j=0}^{n-1}\eta_j\ind_{[t_j,t_{j+1})}(t)
  \]
  dan tulis
  \[
    \Delta_jW:=W(t_{j+1})-W(t_j),
    \qquad
    \Delta_jt:=t_{j+1}-t_j.
  \]
  Maka
  \[
    I^\ast(f)=\sum_{j=0}^{n-1}\eta_j\Delta_jW.
  \]

  Langkah pertama adalah mengkuadratkan bentuk di atas:
  \[
    (I^\ast(f))^2
    =
    \sum_{j=0}^{n-1}\eta_j^2(\Delta_jW)^2
    +
    2\sum_{0\le k<j\le n-1}\eta_j\eta_k\,\Delta_jW\,\Delta_kW.
  \]
  Ambil ekspektasi pada kedua ruas. Kita peroleh
  \begin{equation}
    \label{eq:isometri-awal}
    \E\bigl((I^\ast(f))^2\bigr)
    =
    \sum_{j=0}^{n-1}\E\bigl(\eta_j^2(\Delta_jW)^2\bigr)
    +
    2\sum_{0\le k<j\le n-1}\E\bigl(\eta_j\eta_k\Delta_jW\Delta_kW\bigr).
  \end{equation}

  Sekarang kita hitung suku diagonal. Karena $\eta_j$ adalah
  $\F_{t_j}$-terukur dan $\Delta_jW$ independen terhadap $\F_{t_j}$, maka
  dengan sifat menara dan sifat mengeluarkan faktor terukur diperoleh
  \begin{align*}
    \E\bigl(\eta_j^2(\Delta_jW)^2\bigr)
     & =
    \E\left(
    \E\bigl(\eta_j^2(\Delta_jW)^2\mid\F_{t_j}\bigr)
    \right)                                                \\
     & =
    \E\left(
    \eta_j^2\E\bigl((\Delta_jW)^2\mid\F_{t_j}\bigr)
    \right)                                                \\
     & =
    \E\left(
    \eta_j^2\E\bigl((\Delta_jW)^2\bigr)
    \right)
    \qquad
    \text{(karena $\Delta_jW$ independen dari $\F_{t_j}$)} \\
     & =
    \E(\eta_j^2)\,\Delta_jt.
  \end{align*}

  Berikutnya, kita periksa suku silang. Ambil $k<j$. Karena
  $\eta_k,\eta_j,\Delta_kW$ semuanya $\F_{t_j}$-terukur dan $\Delta_jW$
  independen terhadap $\F_{t_j}$, maka
  \begin{align*}
    \E(\eta_j\eta_k\Delta_jW\Delta_kW)
     & =
    \E\left(
    \E(\eta_j\eta_k\Delta_jW\Delta_kW\mid\F_{t_j})
    \right) \\
     & =
    \E\left(
    \eta_j\eta_k\Delta_kW\,\E(\Delta_jW\mid\F_{t_j})
    \right) \\
     & =
    \E\left(
    \eta_j\eta_k\Delta_kW\,\E(\Delta_jW)
    \right) \\
     & =0,
  \end{align*}
  sebab $\E(\Delta_jW)=0$.

  Dengan hasil ini, semua suku silang pada \eqref{eq:isometri-awal} hilang,
  sehingga
  \[
    \E\bigl((I^\ast(f))^2\bigr)
    =
    \sum_{j=0}^{n-1}\E(\eta_j^2)\Delta_jt.
  \]

  Di sisi lain, karena $f$ konstan pada setiap interval partisi, maka
  \[
    \int_0^\infty f(t)^2\,dt
    =
    \sum_{j=0}^{n-1}\eta_j^2\Delta_jt.
  \]
  Mengambil ekspektasi menghasilkan
  \[
    \E\left(\int_0^\infty f(t)^2\,dt\right)
    =
    \sum_{j=0}^{n-1}\E(\eta_j^2)\Delta_jt.
  \]
  Jadi
  \[
    \E\bigl((I^\ast(f))^2\bigr)
    =
    \E\left(\int_0^\infty f(t)^2\,dt\right).
  \]
  Karena ruas kanan berhingga, maka $I^\ast(f)\in L^2$.
\end{proof}

\begin{lema}
  Untuk setiap $f,g\in M_{\mathrm{step}}^2$ berlaku
  \[
    \E\bigl(I^\ast(f)I^\ast(g)\bigr)
    =
    \E\left(\int_0^\infty f(t)g(t)\,dt\right).
  \]
\end{lema}

\begin{proof}
  Ambil partisi bersama sehingga
  \[
    f(t)=\sum_{j=0}^{n-1}\eta_j\ind_{[t_j,t_{j+1})}(t),
    \qquad
    g(t)=\sum_{j=0}^{n-1}\zeta_j\ind_{[t_j,t_{j+1})}(t).
  \]
  Maka
  \[
    I^\ast(f)I^\ast(g)
    =
    \sum_{j=0}^{n-1}\sum_{k=0}^{n-1}\eta_j\zeta_k\Delta_jW\Delta_kW.
  \]
  Dengan argumen yang sama seperti pada proposisi sebelumnya:
  \begin{enumerate}[label=\arabic*.]
    \item jika $j=k$, maka
          \[
            \E(\eta_j\zeta_j(\Delta_jW)^2)=\E(\eta_j\zeta_j)\Delta_jt;
          \]
    \item jika $j\ne k$, maka suku silang bernilai nol karena salah satu
          inkremen memiliki mean nol dan independen terhadap informasi masa lalu.
  \end{enumerate}
  Akibatnya,
  \[
    \E\bigl(I^\ast(f)I^\ast(g)\bigr)
    =
    \sum_{j=0}^{n-1}\E(\eta_j\zeta_j)\Delta_jt.
  \]
  Sementara itu,
  \[
    \int_0^\infty f(t)g(t)\,dt
    =
    \sum_{j=0}^{n-1}\eta_j\zeta_j\Delta_jt,
  \]
  sehingga setelah diambil ekspektasi diperoleh hasil yang sama.
\end{proof}

\begin{lema}
  Pemetaan $I^\ast:M_{\mathrm{step}}^2\to L^2$ bersifat linear, yaitu
  \[
    I^\ast(\alpha f+\beta g)=\alpha I^\ast(f)+\beta I^\ast(g)
  \]
  untuk setiap $f,g\in M_{\mathrm{step}}^2$ dan $\alpha,\beta\in\R$.
\end{lema}

\begin{proof}
  Tulis $f$ dan $g$ pada partisi yang sama. Maka koefisien dari
  $\alpha f+\beta g$ adalah $\alpha\eta_j+\beta\zeta_j$. Dengan memakai definisi
  integral It\^o untuk proses sederhana,
  \begin{align*}
    I^\ast(\alpha f+\beta g)
     & =
    \sum_{j=0}^{n-1}(\alpha\eta_j+\beta\zeta_j)\Delta_jW \\
     & =
    \alpha\sum_{j=0}^{n-1}\eta_j\Delta_jW
    +
    \beta\sum_{j=0}^{n-1}\zeta_j\Delta_jW                \\
     & =
    \alpha I^\ast(f)+\beta I^\ast(g).
  \end{align*}
\end{proof}

\section{Perluasan Integral It\^o}

\subsection{Ruang M2}

\begin{definisi}
  Ruang $M^2$ adalah himpunan semua proses stokastik $f(t)$ yang memenuhi:
  \begin{enumerate}[label=\arabic*.]
    \item
          \[
            \E\left(\int_0^\infty |f(t)|^2\,dt\right)<\infty;
          \]
    \item terdapat barisan $\{f_n\}\subset M_{\mathrm{step}}^2$ sehingga
          \[
            \lim_{n\to\infty}
            \E\left(\int_0^\infty |f(t)-f_n(t)|^2\,dt\right)=0.
          \]
  \end{enumerate}
\end{definisi}

\begin{contoh}
  Proses deterministik $f(t)=e^{-t}$, $t\ge 0$, termasuk dalam $M^2$.
\end{contoh}

\begin{proof}
  Pertama,
  \[
    \int_0^\infty e^{-2t}\,dt=\frac12,
  \]
  sehingga syarat integrabilitas kuadrat terpenuhi.

  Kedua, untuk setiap $n\in\N$ definisikan aproksimasi tangga
  \[
    f_n(t)
    =
    \sum_{j=0}^{n^2-1}e^{-j/n}\ind_{[j/n,(j+1)/n)}(t).
  \]
  Proses $f_n$ adalah proses sederhana deterministik, maka $f_n\in
    M_{\mathrm{step}}^2$. Karena $e^{-t}$ kontinu dan menurun, diperoleh
  $f_n(t)\to f(t)$ untuk setiap $t\ge 0$. Selain itu,
  \[
    |f_n(t)-f(t)|^2\le 4e^{-2t},
  \]
  dan fungsi pembatas $4e^{-2t}$ terintegralkan pada $[0,\infty)$. Dengan
  teorema konvergensi terdominasi,
  \[
    \int_0^\infty |f_n(t)-f(t)|^2\,dt\to 0.
  \]
  Karena semua fungsi deterministik, mengambil ekspektasi tidak mengubah hasil.
  Jadi $f\in M^2$.
\end{proof}

\subsection{Definisi Integral It\^o pada Ruang M2}

\begin{definisi}
  Jika $f\in M^2$, maka integral It\^o dari $f$ didefinisikan sebagai peubah
  acak $I(f)\in L^2$ yang memenuhi
  \[
    I^\ast(f_n)\xrightarrow{L^2}I(f),
  \]
  untuk setiap barisan $\{f_n\}\subset M_{\mathrm{step}}^2$ yang mendekati $f$
  dalam norma $M^2$.
\end{definisi}

\begin{proposisi}
  Untuk setiap $f\in M^2$, integral stokastik $I(f)$ ada, unik dalam $L^2$, dan
  memenuhi
  \[
    \E\bigl((I(f))^2\bigr)
    =
    \E\left(\int_0^\infty f(t)^2\,dt\right).
  \]
\end{proposisi}

\begin{proof}
  Definisikan norma
  \[
    \|f\|_{M^2}
    :=
    \left(
    \E\left(\int_0^\infty |f(t)|^2\,dt\right)
    \right)^{1/2},
    \qquad
    \|\xi\|_{L^2}:=\bigl(\E(\xi^2)\bigr)^{1/2}.
  \]

  Ambil barisan aproksimasi $\{h_n\}\subset M_{\mathrm{step}}^2$ sehingga
  $\|f-h_n\|_{M^2}\to 0$. Kita tunjukkan bahwa $\{I^\ast(h_n)\}$ adalah barisan
  Cauchy di $L^2$. Dengan linearitas $I^\ast$ dan isometri It\^o untuk proses
  sederhana,
  \begin{align*}
    \|I^\ast(h_n)-I^\ast(h_m)\|_{L^2}^2
     & =
    \|I^\ast(h_n-h_m)\|_{L^2}^2 \\
     & =
    \|h_n-h_m\|_{M^2}^2.
  \end{align*}
  Karena $\{h_n\}$ konvergen dalam $M^2$, barisan tersebut Cauchy. Maka
  $\{I^\ast(h_n)\}$ juga Cauchy di $L^2$.

  Ruang $L^2$ lengkap, sehingga terdapat $Y\in L^2$ dengan
  \[
    I^\ast(h_n)\xrightarrow{L^2}Y.
  \]
  Kita definisikan $I(f):=Y$.

  Berikutnya, kita buktikan bahwa definisi ini tidak bergantung pada pemilihan
  barisan aproksimasi. Misalkan $\{g_n\}\subset M_{\mathrm{step}}^2$ juga
  memenuhi $g_n\to f$ dalam $M^2$. Maka
  \begin{align*}
    \|I^\ast(h_n)-I^\ast(g_n)\|_{L^2}
     & =
    \|I^\ast(h_n-g_n)\|_{L^2} \\
     & =
    \|h_n-g_n\|_{M^2}         \\
     & \le
    \|h_n-f\|_{M^2}+\|g_n-f\|_{M^2}\to 0.
  \end{align*}
  Jadi kedua barisan integral mempunyai limit $L^2$ yang sama. Dengan demikian,
  $I(f)$ terdefinisi dengan baik dan unik.

  Terakhir, dari isometri untuk $h_n$ diperoleh
  \[
    \|I^\ast(h_n)\|_{L^2}=\|h_n\|_{M^2}.
  \]
  Ambil limit ketika $n\to\infty$. Karena norma kontinu terhadap limit $L^2$,
  didapat
  \[
    \|I(f)\|_{L^2}=\|f\|_{M^2}.
  \]
  Dengan mengkuadratkan kedua ruas, hasil yang diinginkan diperoleh:
  \[
    \E\bigl((I(f))^2\bigr)
    =
    \E\left(\int_0^\infty f(t)^2\,dt\right).
  \]
\end{proof}

\begin{lema}
  Jika $f\in M_{\mathrm{step}}^2$, maka integral pada ruang $M^2$ konsisten
  dengan definisi awal, yaitu $I(f)=I^\ast(f)$.
\end{lema}

\begin{proof}
  Ambil sebarang barisan $\{f_n\}\subset M_{\mathrm{step}}^2$ yang mendekati $f$
  dalam $M^2$. Dengan isometri It\^o untuk proses sederhana,
  \[
    \E\bigl(|I^\ast(f)-I^\ast(f_n)|^2\bigr)
    =
    \E\left(\int_0^\infty |f(t)-f_n(t)|^2\,dt\right)\to 0.
  \]
  Jadi $I^\ast(f_n)\xrightarrow{L^2}I^\ast(f)$. Berdasarkan definisi integral
  pada $M^2$, limit ini harus sama dengan $I(f)$. Maka $I(f)=I^\ast(f)$.
\end{proof}

\subsection{Integral pada Interval Berhingga}

\begin{definisi}
  Untuk $T>0$, didefinisikan
  \[
    M_T^2:=\{f:\ind_{[0,T]}f\in M^2\}.
  \]
  Integral It\^o pada interval berhingga didefinisikan oleh
  \[
    I_T(f):=I(\ind_{[0,T]}f)=\int_0^T f(t)\,dW(t).
  \]
\end{definisi}

\begin{lema}
  Jika $f\in M_{\mathrm{step}}^2$, maka proses
  \[
    M_t:=\int_0^t f(s)\,dW(s), \qquad t\ge 0,
  \]
  merupakan martingale.
\end{lema}

\begin{proof}
  Ambil $0\le s<t$ dan perhalus partisi sehingga $s$ dan $t$ menjadi titik
  partisi. Maka ada indeks $k<m$ dengan
  \[
    s=t_k,\qquad t=t_m,
  \]
  dan
  \[
    M_t=\sum_{j=0}^{m-1}\eta_j\Delta_jW,
    \qquad
    M_s=\sum_{j=0}^{k-1}\eta_j\Delta_jW.
  \]
  Hitung ekspektasi bersyarat terhadap $\F_s$:
  \[
    \E(M_t\mid\F_s)
    =
    \sum_{j=0}^{m-1}\E(\eta_j\Delta_jW\mid\F_s).
  \]
  Untuk $j<k$, peubah $\eta_j\Delta_jW$ sudah $\F_s$-terukur, sehingga
  \[
    \E(\eta_j\Delta_jW\mid\F_s)=\eta_j\Delta_jW.
  \]
  Untuk $j\ge k$, karena $\F_s\subseteq \F_{t_j}$, sifat menara memberi
  \[
    \E(\eta_j\Delta_jW\mid\F_s)
    =
    \E\bigl(\E(\eta_j\Delta_jW\mid\F_{t_j})\mid\F_s\bigr).
  \]
  Selanjutnya $\eta_j$ adalah $\F_{t_j}$-terukur dan $\Delta_jW$ independen
  terhadap $\F_{t_j}$ dengan mean nol, maka
  \[
    \E(\eta_j\Delta_jW\mid\F_{t_j})
    =
    \eta_j\E(\Delta_jW\mid\F_{t_j})
    =
    \eta_j\E(\Delta_jW)=0.
  \]
  Jadi semua suku dengan $j\ge k$ hilang, sehingga
  \[
    \E(M_t\mid\F_s)=\sum_{j=0}^{k-1}\eta_j\Delta_jW=M_s.
  \]
  Dengan demikian, $\{M_t\}_{t\ge 0}$ adalah martingale.
\end{proof}

\section{Kriteria Eksistensi Integral It\^o}

\subsection{Proses Teradaptasi dengan Lintasan Kontinu}

\begin{teorema}
  Misalkan $f(t)$ proses stokastik yang teradaptasi terhadap filtrasi
  $\{\F_t\}_{t\ge 0}$ dan memiliki lintasan kontinu hampir pasti. Maka:
  \begin{enumerate}[label=\arabic*.]
    \item jika
          \[
            \E\left(\int_0^\infty |f(t)|^2\,dt\right)<\infty,
          \]
          maka $f\in M^2$;
    \item jika untuk suatu $T>0$ berlaku
          \[
            \E\left(\int_0^T |f(t)|^2\,dt\right)<\infty,
          \]
          maka $f\in M_T^2$.
  \end{enumerate}
\end{teorema}

\begin{proof}
  Kita buktikan butir pertama terlebih dahulu. Untuk setiap $n\in\N$,
  konstruksikan proses sederhana acak
  \[
    f_n(t)=
    \begin{cases}
      n\displaystyle\int_{\frac{k-1}{n}}^{\frac{k}{n}} f(s)\,ds,
         & \text{jika }\frac{k}{n}\le t<\frac{k+1}{n},\ k=1,2,\ldots,n^2-1, \\[3mm]
      0, & \text{jika }t\ge n \text{ atau } 0\le t<\frac1n.
    \end{cases}
  \]
  Koefisien pada interval $\left[\frac{k}{n},\frac{k+1}{n}\right)$ adalah
  \[
    \eta_{k,n}:=
    n\int_{\frac{k-1}{n}}^{\frac{k}{n}} f(s)\,ds.
  \]
  Karena $f(s)$ teradaptasi, maka untuk $s\le \frac{k}{n}$ peubah $f(s)$
  terukur terhadap $\F_{k/n}$. Integralnya juga $\F_{k/n}$-terukur. Jadi
  $\eta_{k,n}$ adalah koefisien yang boleh dipakai pada interval itu.

  Selanjutnya, kita periksa bahwa $f_n$ terintegralkan kuadrat. Dengan
  ketaksamaan Cauchy--Schwarz,
  \[
    \left|
    \int_{\frac{k-1}{n}}^{\frac{k}{n}}f(s)\,ds
    \right|^2
    \le
    \frac1n\int_{\frac{k-1}{n}}^{\frac{k}{n}}|f(s)|^2\,ds.
  \]
  Kalikan kedua ruas dengan $n^2$, lalu integralkan terhadap $t$ pada interval
  yang panjangnya $\frac1n$. Kita memperoleh
  \[
    \int_{\frac{k}{n}}^{\frac{k+1}{n}}|f_n(t)|^2\,dt
    \le
    \int_{\frac{k-1}{n}}^{\frac{k}{n}}|f(s)|^2\,ds.
  \]
  Menjumlahkan untuk semua $k$ memberi
  \[
    \int_0^\infty |f_n(t)|^2\,dt
    \le
    \int_0^\infty |f(t)|^2\,dt
    \qquad \text{hampir pasti.}
  \]
  Setelah diambil ekspektasi, ruas kanan berhingga menurut hipotesis, sehingga
  $f_n\in M_{\mathrm{step}}^2$.

  Berikutnya, kita buktikan bahwa $f_n\to f$ dalam norma $M^2$. Ambil $N>0$.
  Pecah integral galat menjadi bagian kompak dan bagian ekor:
  \begin{align*}
    \int_0^\infty |f(t)-f_n(t)|^2\,dt
     & =
    \int_0^N |f(t)-f_n(t)|^2\,dt
    +
    \int_N^\infty |f(t)-f_n(t)|^2\,dt \\
     & \le
    \int_0^N |f(t)-f_n(t)|^2\,dt
    +
    2\int_N^\infty |f(t)|^2\,dt
    +
    2\int_N^\infty |f_n(t)|^2\,dt.
  \end{align*}
  Dari estimasi sebelumnya, integral ketiga dibatasi oleh
  \[
    \int_N^\infty |f_n(t)|^2\,dt
    \le
    \int_{N-1}^\infty |f(t)|^2\,dt.
  \]
  Jadi
  \begin{equation}
    \label{eq:tail-bound}
    \int_0^\infty |f(t)-f_n(t)|^2\,dt
    \le
    \int_0^N |f(t)-f_n(t)|^2\,dt
    +
    4\int_{N-1}^\infty |f(t)|^2\,dt.
  \end{equation}

  Karena lintasan $f$ kontinu hampir pasti, maka pada setiap interval kompak
  $[0,N]$ lintasannya seragam kontinu hampir pasti. Rataan lokal pada definisi
  $f_n$ karena itu konvergen ke nilai fungsi, sehingga
  \[
    \int_0^N |f(t)-f_n(t)|^2\,dt\to 0
    \qquad \text{hampir pasti.}
  \]
  Sementara itu, karena $\int_0^\infty |f(t)|^2\,dt<\infty$ hampir pasti, maka
  \[
    \int_{N-1}^\infty |f(t)|^2\,dt\to 0
    \qquad \text{ketika }N\to\infty.
  \]
  Dengan \eqref{eq:tail-bound} kita peroleh
  \[
    \int_0^\infty |f(t)-f_n(t)|^2\,dt\to 0
    \qquad \text{hampir pasti.}
  \]

  Selain itu,
  \[
    |f(t)-f_n(t)|^2\le 2|f(t)|^2+2|f_n(t)|^2\le 4|f(t)|^2,
  \]
  dan $4\int_0^\infty |f(t)|^2\,dt$ terintegralkan. Maka, dengan teorema
  konvergensi terdominasi,
  \[
    \E\left(\int_0^\infty |f(t)-f_n(t)|^2\,dt\right)\to 0.
  \]
  Jadi $f\in M^2$.

  Untuk butir kedua, cukup terapkan butir pertama pada proses
  $\ind_{[0,T]}f$. Proses ini tetap teradaptasi dan memenuhi
  \[
    \E\left(\int_0^\infty |\ind_{[0,T]}(t)f(t)|^2\,dt\right)
    =
    \E\left(\int_0^T |f(t)|^2\,dt\right)<\infty.
  \]
  Jadi $\ind_{[0,T]}f\in M^2$, atau ekuivalen dengan $f\in M_T^2$.
\end{proof}

\subsection{Contoh Keanggotaan pada Ruang MT2}

\begin{lema}
  Proses Wiener $W(t)$ termasuk dalam $M_T^2$ untuk setiap $T>0$.
\end{lema}

\begin{proof}
  Proses Wiener jelas teradaptasi terhadap filtrasi alaminya dan memiliki
  lintasan kontinu hampir pasti. Maka, berdasarkan teorema sebelumnya, cukup
  diperiksa bahwa
  \[
    \E\left(\int_0^T W(t)^2\,dt\right)<\infty.
  \]
  Dengan teorema Fubini--Tonelli,
  \[
    \E\left(\int_0^T W(t)^2\,dt\right)
    =
    \int_0^T \E(W(t)^2)\,dt.
  \]
  Karena $W(t)\sim\mathcal{N}(0,t)$, maka $\E(W(t)^2)=t$. Akibatnya,
  \[
    \E\left(\int_0^T W(t)^2\,dt\right)
    =
    \int_0^T t\,dt
    =
    \frac{T^2}{2}
    <\infty.
  \]
  Jadi $W\in M_T^2$.
\end{proof}

\begin{lema}
  Proses $W(t)^2$ termasuk dalam $M_T^2$ untuk setiap $T>0$.
\end{lema}

\begin{proof}
  Karena $W(t)$ adalah $\F_t$-terukur, maka $W(t)^2$ juga $\F_t$-terukur. Karena
  lintasan $W(t)$ kontinu hampir pasti, maka lintasan $W(t)^2$ juga kontinu
  hampir pasti.

  Jadi, berdasarkan teorema kriteria eksistensi, cukup dibuktikan bahwa
  \[
    \E\left(\int_0^T |W(t)^2|^2\,dt\right)
    =
    \E\left(\int_0^T W(t)^4\,dt\right)<\infty.
  \]
  Dengan Fubini--Tonelli,
  \[
    \E\left(\int_0^T W(t)^4\,dt\right)
    =
    \int_0^T \E(W(t)^4)\,dt.
  \]
  Karena $W(t)\sim\mathcal{N}(0,t)$, momen keempat peubah normal bernilai
  \[
    \E(W(t)^4)=3t^2.
  \]
  Maka
  \[
    \E\left(\int_0^T W(t)^4\,dt\right)
    =
    \int_0^T 3t^2\,dt
    =
    T^3
    <\infty.
  \]
  Jadi $W(t)^2\in M_T^2$.
\end{proof}

\section{Proses Progresif Terukur}

\begin{definisi}
  Proses stokastik $f(t,\omega)$ disebut \textbf{progresif terukur} jika untuk
  setiap $T>0$, pemetaan
  \[
    (t,\omega)\longmapsto f(t,\omega),
    \qquad 0\le t\le T,
  \]
  terukur terhadap $\B([0,T])\otimes\F_T$.
\end{definisi}

\begin{contoh}
  Proses Wiener $\{W(t)\}_{t\ge 0}$ adalah proses progresif terukur.
\end{contoh}

\begin{proof}
  Ambil $T>0$. Untuk hampir setiap $\omega$, lintasan
  $t\mapsto W(t,\omega)$ kontinu pada $[0,T]$, sehingga terukur terhadap
  $\B([0,T])$ sebagai fungsi dari $t$. Selain itu, untuk setiap $t\le T$, peubah
  acak $W(t)$ adalah $\F_t$-terukur, dan karena $\F_t\subseteq \F_T$, maka
  $W(t)$ juga $\F_T$-terukur.

  Jadi pemetaan $(t,\omega)\mapsto W(t,\omega)$ terukur terhadap
  $\B([0,T])\otimes\F_T$. Karena $T$ dipilih sebarang, proses Wiener progresif
  terukur.
\end{proof}

\begin{teorema}
  Berlaku karakterisasi berikut.
  \begin{enumerate}[label=\arabic*.]
    \item Ruang $M^2$ terdiri dari semua proses progresif terukur yang memenuhi
          \[
            \E\left(\int_0^\infty |f(t)|^2\,dt\right)<\infty.
          \]
    \item Ruang $M_T^2$ terdiri dari semua proses progresif terukur yang memenuhi
          \[
            \E\left(\int_0^T |f(t)|^2\,dt\right)<\infty.
          \]
  \end{enumerate}
\end{teorema}

\begin{proof}
  Pembuktian untuk $M_T^2$ serupa, sehingga cukup dibahas $M^2$.

  $(\Rightarrow)$ Misalkan $f\in M^2$. Menurut definisi, ada barisan
  $\{f_n\}\subset M_{\mathrm{step}}^2$ yang mendekati $f$ dalam norma $M^2$.
  Setiap proses sederhana acak adalah progresif terukur karena pada tiap
  interval partisi ia berbentuk hasil kali fungsi indikator waktu dengan
  koefisien yang $\F_{t_j}$-terukur. Dari konvergensi $L^2$ pada ruang produk,
  dapat diambil subsekuen yang konvergen hampir di mana-mana; limit hampir di
  mana-mana dari fungsi progresif terukur tetap progresif terukur. Jadi $f$
  progresif terukur. Syarat integrabilitas kuadrat sudah terdapat dalam definisi
  $M^2$.

  $(\Leftarrow)$ Misalkan $f$ progresif terukur dan
  \[
    \E\left(\int_0^\infty |f(t)|^2\,dt\right)<\infty.
  \]
  Secara umum, fungsi terukur pada ruang produk dapat didekati dalam norma
  $L^2$ oleh fungsi sederhana terukur. Karena $f$ progresif terukur, aproksimasi
  sederhana dapat dipilih konsisten dengan struktur progresif, sehingga
  menghasilkan barisan proses sederhana acak $\{f_n\}\subset M_{\mathrm{step}}^2$
  dengan
  \[
    \E\left(\int_0^\infty |f(t)-f_n(t)|^2\,dt\right)\to 0.
  \]
  Dengan demikian, $f\in M^2$.
\end{proof}

\section{Contoh Perhitungan Integral It\^o}

\subsection{Verifikasi Integral It\^o untuk Proses Wiener}

\begin{proposisi}
  Untuk setiap $T>0$ berlaku
  \[
    \int_0^T W(t)\,dW(t)
    =
    \frac12 W(T)^2-\frac12 T.
  \]
\end{proposisi}

\begin{proof}
  Menurut lema sebelumnya, $W(t)\in M_T^2$, sehingga integral di atas
  terdefinisi. Ambil partisi seragam
  \[
    t_j=\frac{jT}{n},
    \qquad
    j=0,1,\ldots,n.
  \]
  Definisikan proses sederhana
  \[
    h_n(t)=\sum_{j=0}^{n-1}W(t_j)\ind_{[t_j,t_{j+1})}(t).
  \]
  Maka $h_n\to W$ dalam $M_T^2$, dan
  \[
    I^\ast(h_n)=\sum_{j=0}^{n-1}W(t_j)\Delta_jW.
  \]

  Gunakan identitas aljabar
  \[
    a(b-a)=\frac12\bigl(b^2-a^2-(b-a)^2\bigr).
  \]
  Dengan $a=W(t_j)$ dan $b=W(t_{j+1})$, kita peroleh
  \[
    W(t_j)\Delta_jW
    =
    \frac12\left(
    W(t_{j+1})^2-W(t_j)^2-(\Delta_jW)^2
    \right).
  \]
  Jumlahkan dari $j=0$ sampai $n-1$:
  \[
    I^\ast(h_n)
    =
    \frac12\sum_{j=0}^{n-1}\bigl(W(t_{j+1})^2-W(t_j)^2\bigr)
    -
    \frac12\sum_{j=0}^{n-1}(\Delta_jW)^2.
  \]

  Suku pertama adalah jumlah teleskopik, sehingga
  \[
    \sum_{j=0}^{n-1}\bigl(W(t_{j+1})^2-W(t_j)^2\bigr)
    =
    W(T)^2-W(0)^2
    =
    W(T)^2.
  \]
  Jadi
  \[
    I^\ast(h_n)
    =
    \frac12W(T)^2-\frac12\sum_{j=0}^{n-1}(\Delta_jW)^2.
  \]

  Sekarang dipakai fakta variasi kuadrat proses Wiener:
  \[
    \sum_{j=0}^{n-1}(\Delta_jW)^2\xrightarrow{L^2}T.
  \]
  Karena $I^\ast(h_n)\xrightarrow{L^2}\int_0^T W(t)\,dW(t)$, maka dengan
  mengambil limit $L^2$ diperoleh
  \[
    \int_0^T W(t)\,dW(t)
    =
    \frac12W(T)^2-\frac12T.
  \]
\end{proof}

\subsection{Verifikasi Integral It\^o untuk Fungsi Deterministik t}

\begin{proposisi}
  Untuk setiap $T>0$ berlaku
  \[
    \int_0^T t\,dW(t)=TW(T)-\int_0^T W(t)\,dt.
  \]
\end{proposisi}

\begin{proof}
  Ambil partisi seragam $t_j=\frac{jT}{n}$ dan definisikan
  \[
    h_n(t)=\sum_{j=0}^{n-1} t_j\ind_{[t_j,t_{j+1})}(t).
  \]
  Maka
  \[
    I^\ast(h_n)=\sum_{j=0}^{n-1} t_j\Delta_jW.
  \]
  Gunakan identitas
  \[
    c(b-a)=(db-ca)-b(d-c).
  \]
  Dengan $c=t_j$, $d=t_{j+1}$, $a=W(t_j)$, dan $b=W(t_{j+1})$, diperoleh
  \[
    t_j\Delta_jW
    =
    \bigl(t_{j+1}W(t_{j+1})-t_jW(t_j)\bigr)-W(t_{j+1})\Delta_jt.
  \]
  Jumlahkan seluruh suku:
  \[
    I^\ast(h_n)
    =
    \sum_{j=0}^{n-1}\bigl(t_{j+1}W(t_{j+1})-t_jW(t_j)\bigr)
    -
    \sum_{j=0}^{n-1}W(t_{j+1})\Delta_jt.
  \]
  Jumlah pertama teleskopik, sehingga
  \[
    I^\ast(h_n)=TW(T)-\sum_{j=0}^{n-1}W(t_{j+1})\Delta_jt.
  \]
  Karena lintasan Wiener kontinu hampir pasti, jumlah Riemann di ruas kanan
  konvergen ke integral Riemann biasa:
  \[
    \sum_{j=0}^{n-1}W(t_{j+1})\Delta_jt\to \int_0^T W(t)\,dt.
  \]
  Ambil limit dan gunakan $h_n\to t$ dalam $M_T^2$ untuk memperoleh
  \[
    \int_0^T t\,dW(t)=TW(T)-\int_0^T W(t)\,dt.
  \]
\end{proof}

\subsection{Verifikasi Integral It\^o untuk Kuadrat Proses Wiener}

\begin{proposisi}
  Untuk setiap $T>0$ berlaku
  \[
    \int_0^T W(t)^2\,dW(t)
    =
    \frac13 W(T)^3-\int_0^T W(t)\,dt.
  \]
\end{proposisi}

\begin{proof}
  Dari hasil sebelumnya, $W(t)^2\in M_T^2$. Ambil partisi
  \[
    0=t_0<t_1<\cdots<t_n=T,
  \]
  dan tulis $\Delta_jW=W(t_{j+1})-W(t_j)$ serta $\Delta_jt=t_{j+1}-t_j$. Untuk
  aproksimasi kiri diperoleh
  \[
    I^\ast(h_n)=\sum_{j=0}^{n-1}W(t_j)^2\Delta_jW.
  \]

  Gunakan identitas aljabar
  \[
    a^2(b-a)
    =
    \frac13(b^3-a^3)-a(b-a)^2-\frac13(b-a)^3.
  \]
  Substitusi $a=W(t_j)$ dan $b=W(t_{j+1})$ memberi
  \[
    W(t_j)^2\Delta_jW
    =
    \frac13\bigl(W(t_{j+1})^3-W(t_j)^3\bigr)
    -
    W(t_j)(\Delta_jW)^2
    -
    \frac13(\Delta_jW)^3.
  \]
  Setelah dijumlahkan, kita peroleh
  \[
    I^\ast(h_n)
    =
    \frac13 W(T)^3
    -
    \sum_{j=0}^{n-1}W(t_j)(\Delta_jW)^2
    -
    \frac13\sum_{j=0}^{n-1}(\Delta_jW)^3.
  \]

  Sekarang kita analisis dua jumlah terakhir.

  \textbf{Langkah 1:} Tunjukkan bahwa
  \[
    \sum_{j=0}^{n-1}(\Delta_jW)^3\xrightarrow{L^2}0.
  \]
  Misalkan
  \[
    S_n:=\sum_{j=0}^{n-1}(\Delta_jW)^3.
  \]
  Karena $\Delta_jW\sim\mathcal{N}(0,\Delta_jt)$, maka
  \[
    \E((\Delta_jW)^3)=0,
    \qquad
    \E((\Delta_jW)^6)=15(\Delta_jt)^3.
  \]
  Dengan independensi inkremen Wiener, semua suku silang dalam $\E(S_n^2)$
  bernilai nol, sehingga
  \[
    \E(S_n^2)=15\sum_{j=0}^{n-1}(\Delta_jt)^3.
  \]
  Jika $\|\Pi\|:=\max_j\Delta_jt$, maka
  \[
    \sum_{j=0}^{n-1}(\Delta_jt)^3
    \le
    \|\Pi\|^2\sum_{j=0}^{n-1}\Delta_jt
    =
    T\|\Pi\|^2.
  \]
  Jadi
  \[
    \E(S_n^2)\le 15T\|\Pi\|^2\to 0.
  \]
  Ini membuktikan $S_n\to 0$ dalam $L^2$.

  \textbf{Langkah 2:} Tunjukkan bahwa
  \[
    \sum_{j=0}^{n-1}W(t_j)(\Delta_jW)^2
    \xrightarrow{L^2}
    \int_0^T W(t)\,dt.
  \]
  Pecah jumlah tersebut menjadi
  \[
    \sum_{j=0}^{n-1}W(t_j)(\Delta_jW)^2
    =
    \sum_{j=0}^{n-1}W(t_j)\bigl((\Delta_jW)^2-\Delta_jt\bigr)
    +
    \sum_{j=0}^{n-1}W(t_j)\Delta_jt.
  \]

  Jumlah kedua adalah jumlah Riemann dari fungsi kontinu $W(t)$, sehingga
  \[
    \sum_{j=0}^{n-1}W(t_j)\Delta_jt
    \to
    \int_0^T W(t)\,dt
  \]
  hampir pasti.

  Untuk jumlah pertama, misalkan
  \[
    R_n:=\sum_{j=0}^{n-1}W(t_j)\bigl((\Delta_jW)^2-\Delta_jt\bigr).
  \]
  Kita hitung momen keduanya. Karena $W(t_j)$ adalah $\F_{t_j}$-terukur dan
  $\Delta_jW$ independen terhadap $\F_{t_j}$, serta
  \[
    \E\bigl((\Delta_jW)^2-\Delta_jt\bigr)=0,
  \]
  suku silang kembali hilang. Maka
  \begin{align*}
    \E(R_n^2)
     & =
    \sum_{j=0}^{n-1}
    \E\left(
    W(t_j)^2\bigl((\Delta_jW)^2-\Delta_jt\bigr)^2
    \right) \\
     & =
    \sum_{j=0}^{n-1}
    \E\bigl(W(t_j)^2\bigr)\,
    \E\left(\bigl((\Delta_jW)^2-\Delta_jt\bigr)^2\right).
  \end{align*}
  Karena $\E(W(t_j)^2)=t_j$ dan untuk peubah normal terpusat berlaku
  \[
    \E\left(\bigl((\Delta_jW)^2-\Delta_jt\bigr)^2\right)=2(\Delta_jt)^2,
  \]
  diperoleh
  \[
    \E(R_n^2)=2\sum_{j=0}^{n-1}t_j(\Delta_jt)^2.
  \]
  Karena $t_j\le T$, maka
  \[
    \E(R_n^2)
    \le
    2T\sum_{j=0}^{n-1}(\Delta_jt)^2
    \le
    2T^2\|\Pi\|\to 0.
  \]
  Jadi $R_n\to 0$ dalam $L^2$. Akibatnya,
  \[
    \sum_{j=0}^{n-1}W(t_j)(\Delta_jW)^2
    \xrightarrow{L^2}
    \int_0^T W(t)\,dt.
  \]

  Mengambil limit pada rumus untuk $I^\ast(h_n)$ menghasilkan
  \[
    \int_0^T W(t)^2\,dW(t)
    =
    \frac13W(T)^3-\int_0^T W(t)\,dt.
  \]
\end{proof}

\section{Penutup}

\subsection{Kesimpulan}

Berdasarkan pembahasan yang telah disusun, dapat disimpulkan beberapa hal
berikut.
\begin{enumerate}[label=\arabic*.]
  \item Kalkulus It\^o muncul karena kalkulus klasik tidak memadai untuk
        menangani proses Wiener yang lintasannya kontinu tetapi sangat tidak
        teratur.
  \item Integral It\^o mula-mula didefinisikan pada proses sederhana acak dengan
        mengambil koefisien di titik kiri interval partisi.
  \item Sifat terpenting integral It\^o adalah isometri
        \[
          \E\bigl((I(f))^2\bigr)=\E\left(\int f(t)^2\,dt\right),
        \]
        yang menjadi dasar perluasan integral dari $M_{\mathrm{step}}^2$ ke
        ruang $M^2$.
  \item Untuk proses yang teradaptasi, berlintasan kontinu hampir pasti, dan
        memenuhi integrabilitas kuadrat, integral It\^o pasti ada.
  \item Identitas seperti
        \[
          \int_0^T W(t)\,dW(t)=\frac12W(T)^2-\frac12T
        \]
        menunjukkan adanya koreksi tambahan yang tidak muncul dalam kalkulus
        biasa, yaitu akibat variasi kuadrat proses Wiener.
\end{enumerate}

\subsection{Saran}

Sebagai kelanjutan dari laporan ini, materi dapat diperluas ke formula It\^o
multidimensi, persamaan diferensial stokastik, martingale representation
theorem, dan penerapan pada model Black--Scholes maupun proses difusi lain.

\begin{thebibliography}{9}

  \bibitem{karatzas}
  Ioannis Karatzas dan Steven E. Shreve,
  \textit{Brownian Motion and Stochastic Calculus},
  2nd ed., Springer, 1991.

  \bibitem{oksendal}
  Bernt {\O}ksendal,
  \textit{Stochastic Differential Equations: An Introduction with Applications},
  6th ed., Springer, 2003.

  \bibitem{revuz}
  Daniel Revuz dan Marc Yor,
  \textit{Continuous Martingales and Brownian Motion},
  3rd ed., Springer, 1999.

\end{thebibliography}

\end{document}
